% !PRINTED 14
\weberpage{14}

%% [running head, not body text -- commented out 2026-08-07] Page furniture (running
%% head, folio, signature) is excluded from this edition by standing policy.
%% \gesperrt{Erster Abschnitt.} \hfill 14 \hfill \gesperrt{\S.\,4.}

But on account of the property of normal divisors, $N$ is permutable with every $b$, hence $b_h N = N b_h$, and consequently
\[
N_h N_k = N N b_h b_k.
\]

Because $b_h b_k$ is contained in $P$, $N b_h b_k$ is identical with one of the subsidiary groups (5), and since $NN = N$ may be set (by 1.), it follows, if we set $N b_h b_k = N_l$,
\weq{6}{N_h N_k = N_l.}

With this, the property \S.\,1, 1. of groups is established for the composition of the subsidiary groups. That \S.\,1, 2. also holds, namely the associative law, we have already pointed out. It therefore remains to prove \S.\,1, 3., namely that, if
\weq{7}{N_h N_i = N_k N_l}
is set, and $N_h = N_k$, then also $N_i = N_l$, and if $N_i = N_l$, then also $N_h = N_k$ must hold.

By the representation (5), however, we can represent (7) in the twofold manner:
\weq{8}{\begin{aligned}
N b_h b_i &= N b_k b_l \\
b_h b_i N &= b_k b_l N.
\end{aligned}}

If now $N_i = N_l$, we may also assume $b_i = b_l$, and obtain from (8), by composition with $b_i^{-1}$,
\[
N b_h = N b_k;
\]
and if $N b_h = N b_k$, we assume $b_h = b_k$ and obtain from (8), by composition with $b_h^{-1}$,
\[
b_i N = b_l N,
\]
hence also $N_i = N_l$.

It is thus proved that the system of subsidiary groups becomes a group under the composition of the parts; that in this group the normal divisor $N$ is the unit element follows immediately from 1., because by it
\[
N N_k = N N b_k = N_k
\]
holds. And because $N b_k N b_k^{-1} = N b_k b_k^{-1} = N$, $N b_k$ and $N b_k^{-1}$ are opposite elements.

The group of the quantities $N_k$, whose existence is hereby proved, we call the \gesperrt{group of the subsidiary groups}, or the \gesperrt{group complementary to $N$} with respect to $P$. We
%CONT>>
